\providecommand\ExamFontPath{../../../字体/}
\documentclass[fontset=exam]{exam-zh}
\usepackage{longtable,booktabs,array,calc,multirow,fvextra}
\usepackage[normalem]{ulem}
\providecommand{\hl}[1]{#1}
\providecommand{\ul}[1]{\uline{#1}}
\DefineVerbatimEnvironment{verbatim}{Verbatim}{breaklines=true,fontsize=\small}
\providecommand{\tightlist}{\setlength{\itemsep}{0pt}\setlength{\parskip}{0pt}}
\newcommand{\pandocbounded}[1]{#1}
\setlength{\emergencystretch}{3em}
\setlength{\tabcolsep}{2pt}
\examsetup{page/size=a4paper}
\title{UNIX 期中试卷}
\subject[9em]{UNIX 系统编程}
\setcounter{tocdepth}{1}
\begin{document}
\maketitle
\fancyfoot[C]{第 \thepage 页（共 \pageref*{LastPage} 页）}
\pdfbookmark[1]{一、}{section1}

\par\medskip{\heiti

一、

\par}\smallskip

1、如果两个常规文件的长度一样，那么它们是否占用相同大小的磁盘空间？为什么？

2、一个进程可以给其他进程中指定的线程发送信号吗？线程可以拥有自己的信号处理程序吗？线程可以有自己的信号屏蔽码吗？

3、使用系统调用 read 和 write 读/写磁盘文件时，是否一定会访问磁盘驱动器？为什么？

4、如果超级用户的一个文件的设置权限是，文件的所有者可以读写，但是同组及其他用户均不可以读写，那么普通进程是否就不能读写这个文件？为什么？

5、当创建一个进程时，其实际用户 ID 与有效用户 ID 是否总是一致？如果不是，那么在什么情况下它们是不同的？

6、在作业控制中，进程组和进程对话有什么用途？

7、如何在不执行相关的信号处理函数的前提下清除一个未决信号？

8、为什么要避免进程僵尸？请给出两种避免进程僵尸的方法？

9、当分别调用函数 exit 和 \_exit 终止进程时，它们会关闭打开的文件吗？它们会关闭打开的标准 I/O 流吗？

10、什么是系统调用？是否有在 C 语言标准里定义的库函数是系统调用？

11、在进程中删除一个文件后，能否打开这个文件？如果文件已经打开，能否继续访问这个文件？

12、子进程继承父进程打开的文件吗？继承父进程对信号的处理方式吗？继承父进程的信号吗？

13、在操作系统编程环境中，程序中需要互斥的代码叫什么？如何识别程序中需要互斥的代码？

14、如何判断一个函数是否是不可重入的？

15、如果一个进程有 5 个线程，那么它创建的子进程最多有几个线程？最少呢？

\pdfbookmark[1]{二、程序阅读：进程创建与执行}{section18}

\par\medskip{\heiti

二、程序阅读：进程创建与执行（教材第 9～10 页）

\par}\smallskip

1、子进程的进程 ID 等于 0 吗？

2、子进程执行的程序能否带参数？

3、程序语句 exit(127) 在什么情况下被执行？子进程正常结束时的返回值一定是 127 吗？

4、程序中语句 waitpid 起什么作用？如果不用，则可能造成什么后果？

\begin{verbatim}
#include "apue.h"
#include<sys/wait.h>
int main(void)
{
    char buf [MAXLINE];
    pid_t pid;
    int status;
    printf("%% ");
    while(fgets(buf, MAXLINE, stdin)!=NULL) {
        if(buf[strlen(buf)-1] == '\n')
            buf[strlen(buf)-1] = 0;/*replace newline with null*/
        if((pid = fork()) < 0) {
            err_sys("fork error");
        } else if(pid == 0) {
            execlp(buf, buf, (char *)0);
            err_ret("couldn't execute: %s", buf);
            exit(127);
        }
        /*parent*/
        if((pid = waitpid(pid, &status, 0)) < 0)
            err_sys("waitpid error");
        printf("%% ");
    }
    exit(0);
}
\end{verbatim}

\pdfbookmark[1]{三、程序阅读：目录树遍历}{section49}

\par\medskip{\heiti

三、程序阅读：目录树遍历（dopath 函数）

\par}\smallskip

1、dopath 函数返回值的有什么意义？

2、程序遍历目录树的次序

3、在 dopath 函数中，语句（不含引号）``ptr{[}-1{]}=0''起什么作用？是否可以省略？

\end{document}
